\documentclass[../../../main.tex]{subfiles}

\begin{document}

\section{Partial Derivatives}

Finally! We are discussing about derivatives for multivariable
functions. The derivatives are actually more complicated than
single-variable functions, and in this note we will discuss about
partial derivatives, not directional derivatives, of multivariable
functions.

First, recall that single-variable function $f$, its derivative is
defined as the following.
\[
f'(x) = \lim_{h \to 0} \frac{f(x + h) - f(x)}{h}
\]
Geometrically interpreting this, it can be considered as the slope
of the tangent line at the point. Another way of interpreting that
was instantaneous rate of change that the point.

For multivariable functions, it should be similar. But how do we
determine which tangent line to use? If we think about it, it is
actually not very easy to determine because there are many tangent
lines at a specific point. Also, if we think of it as instantaneous
rate of change, the direction of the change can give different values.
This is where directional derivatives comes in, but in this section,
we will be discussing about a more simple derivative that we can
think of from multivariable functions.

For functions with two variables, we don't really have to care about
direction if we slice a function and consider the derivative based
on that slicing line. That slicing line can be a line parallel to
$x$ or $y$ axis. If we consider them as special derivatives,
we can define partial derivatives as the following.

\begin{definition}{}{}
Given a function $f(x, y)$, its \textbf{partial derivative with
respect to $\mathbf{x}$} at point $(x_0, y_0)$ is defined as the
following.
\[
f_x (x_0, y_0)
= \lim_{h \to 0} \frac{f(x_0 + h, y_0) - f(x_0, y_0)}{h}
\]
\end{definition}

Similarly for partial derivative with respect to $y$, we can write
it as $f_y (x_0, y_0) = \lim_{h \to 0} \frac{f(x_0, y_0 + h) -
f(x_0, y_0)}{h}$. This is much easier because by slicing and thinking
only about that slice, it becomes ordinary derivative of a
single-variable function. Another way of thinking partial derivatives
is we are simply putting the other variable as a constant and finding
the derivative with the remaining single-variable function.
Here is an example.

\begin{problem}{}{}
Given $f(x, y) = x^2 + 2x + y^2$, Find $f_x$ and $f_y$ at $(1, 2)$.
\end{problem}
\begin{solution}
Consider the following equation for $f_x(1, 2)$.
\begin{align*}
    f_x (1, 2)
    &= \lim_{h \to 0} \frac{f(1 + h, 2) - f(1, 2)}{h} \\
    &= \lim_{h \to 0} \frac{h^2 + 2h + 1 + 2h + 2 + 4 - 7}{h} \\
    &= \lim_{h \to 0} \frac{h^2 + 4h}{h}
    = \lim_{h \to 0} h + 4
    = 4
\end{align*}
Therefore, $f_x(1, 2) = 4$.

For $f_y(1, 2)$ let's take a different approach. Considering $x$
as a constant and finding the derivative with respect to $y$,
we get $f_y = 2y$. Therefore, $f_y(1, 2) = 4$.
\end{solution}

Before moving on to higher order partial derivatives, here is
another very popular way of writing partial derivatives.
\[
f_x = \frac{\partial f}{\partial x}
\quad\text{and}\quad
f_y = \frac{\partial f}{\partial y}
\]
Now how do we take higher order partial derivative of $f$?

First, notice that $f_x$ and $f_y$ are the first-order partial
derivatives of $f$. Using this, we can find the second order.
\begin{align*}
    f_x &\rightarrow (f_x)_x,\ (f_x)_y \\
    f_y &\rightarrow (f_y)_x,\ (f_y)_y
\end{align*}
Note that we usually suppress the parenthesis and just simply write
$f_{xx}, f_{xy}, f_{yx}, f_{yy}$. Using these, we can also find the
third-order partial derivatives of $f$.
\begin{align*}
    f_{xx} &\rightarrow f_{xxx},\ f_{xxy} \\
    f_{xy} &\rightarrow f_{xyx},\ f_{xyy} \\
    f_{yx} &\rightarrow f_{yxx},\ f_{yxy} \\
    f_{yy} &\rightarrow f_{yyx},\ f_{yyy}
\end{align*}
As long as the resulting partial derivatives are differentiable,
we can continue this process! Notice that if we use the other
notation, then the order in which write the partial derivatives
are flipped. Here is what I mean.
\[
f_{xyy}
= \frac{\partial}{\partial y} \left(
    \frac{\partial}{\partial y} \left(
        \frac{\partial f}{\partial x}
    \right)
\right)
= \frac{\partial^3 f}{\partial y\, \partial y\, \partial x}
\]
With the notation involving the partial symbol, we read it from
right to left. Here is an example.

\begin{problem}{}{}
Given $f(x, y) = x^2y^3 + 2x + 3y^2$, compute $f_{xyy}$.
\end{problem}
\begin{solution}
First, notice that we need to find the partial derivative of $f$
with respect to $x$.
\[
\frac{\partial f}{\partial x}
= 2xy^3 + 2
\]
Continuing, we can find the following.
\[
\frac{\partial^2}{\partial y^2} (2xy^3 + 2)
= \frac{\partial}{\partial y} 6xy^2
= 12xy
\]
Therefore, $\frac{\partial^3 f}{\partial y^2 \partial x} = 12xy$.
\end{solution}

Continuing, we can consider a very important theorem.

\begin{theorem}{Clairaut's Theorem}{Clairaut's Theorem}
For a function $f$ of two variables, if $f_{xy}$ and $f_{yx}$ are
continuous on an open set $A$, then $f_{xy} = f_{yx}$ for all points
in $A$.
\end{theorem}
\begin{proof}
The proof can be done by applying MVT two times for two different
cases. First, define $\Delta h = [f(a + h, b + h) - f(a + h, b)]
- [f(a, b + h) - f(a, b)]$.

For the first case, define $g(x) = f(x, b + h) - f(x, b)$.
By construction, it is evident that $g$ is a single-variable function
and $\Delta h = g(a + h) - g(a)$. By MVT, there exists $c \in
(a, a + h)$ such that $g'(c) = \frac{g(a + h) - g(a)}{a + h - a}$,
or $h g'(c) = g(a + h) - g(a)$. In other words, $\Delta h = hg'(c)$
for some constant $h$ and $c$. Continuing, notice that $g'(c)
= f_x f(c, b + h) - f_x f(c, b)$ by definition. Using MVT again,
there exists $d \in (b, b + h)$ such that $h f_{xy} (c, d) = g'(c)$
since $c$ is a constant. In other words, $\Delta h$ can be written
as $h^2 f_{xy}$ and the following equation holds.
\[
\lim_{h \to 0} \frac{\Delta h}{h^2}
= \lim_{(c, d) \to (a, b)} f_{xy} (c, d)
= f_{xy} (a, b)
\]
For the second case, let $g(y) = f(a + h, y) - f(a, y)$. Notice that
$\Delta h = [f(a + h, b + h) - f(a, b + h)] - [f(a + h, b) - f(a, b)]
= g(b + h) - g(b)$. Similarly by MVT, there exists $c \in (b, b + h)$
such that $h g'(c) = g(b + h) - g(b) = \Delta h$. Moreover, because
$g'(c) = f_y (a + h, c) - f_y (a, c)$, we MVT can be used again to
show that there exists $d \in (a, a + h)$ such that $h f_{yx} (d, c) =
g'(c)$. Therefore, $\Delta h$ can be rewritten as $h^2 f_{yx} (d, c)$,
and the following equation holds.
\[
\lim_{h \to 0} \frac{\Delta h}{h^2}
= \lim_{(d, c) \to (a, b)} f_{yx} (d, c)
= f_{yx} (a, b)
\]
Therefore, $f_{xy} (a, b) = f_{yx} (a, b)$ for all points $(a, b)$
in $\dom(f)$ and the theorem holds.
\end{proof}

In fact, this theorem is not limited to functions of two variables
at all! We can simply let $g(x_i, x_j) = f(x_1, \ldots, x_i, \ldots,
x_j, \ldots, x_n)$ and notice that $g_{x_i x_j} = f_{x_i x_j}$ and
$g_{x_j x_i} = f_{x_j x_i}$. Because $g_{x_i x_j} = g_{x_j x_i}$,
$f_{x_i x_j} = f_{x_j x_i}$ for function $f$ of $n$ variables where
$n > 2$.

Notice that this theorem implies that we can switch the order of
writing $x,y$ for higher order partial derivatives. Meaning,
all it matters is the number of $x$ and $y$ in the subscript,
not their orders when taking higher order partials.

\begin{important}
Finally, for a function $f$ of $n$ variables, our discussion can
be generalized as the following.
\[
f_{x_i} (x_1, \ldots, x_n)
= \lim_{h \to 0} \frac{
    f(x_1, \ldots, x_i + h, \ldots, x_n - f(x_1, \ldots, x_n)
}{h}
\]
\end{important}



\end{document}


